\section*{Hoofdstuk \ref{ch:b3:homogeneous-catalysis} --- Homogene katalyse in de industrie}

\begin{solution}{exo:b3:homogeneous-catalysis:1}
$\ce{RhCl(PPh3)3}$: 16, Rh(I) $d^8$. $\ce{RhCl(PPh3)2}$: 14, Rh(I). $\ce{RhClH2(PPh3)2}$: 16, Rh(III) $d^6$.
Alkeencomplex: 18, Rh(III). Alkylhydride: 16, Rh(III). De reductieve
eliminatie keert terug naar 14, Rh(I).
\end{solution}

\begin{solution}{exo:b3:homogeneous-catalysis:2}
Butanal, $\ce{CH3CH2CH2CHO}$ (onvertakt), en 2-methylpropanal, $\ce{(CH3)2CHCHO}$ (vertakt).
\end{solution}

\begin{solution}{exo:b3:homogeneous-catalysis:3}
Diethylcyclopent-3-een-1,1-dicarboxylaat, een vijfring, onder verlies van
etheen: de twee eindstandige groepen $\ce{CH2}$ vertrekken als $\ce{C2H4}$ en de twee inwendige CH-koolstofatomen
worden verbonden.
\end{solution}

\begin{solution}{exo:b3:homogeneous-catalysis:4}
Methyl-(\textit{E})-3-fenylprop-2-enoaat (methylcinnamaat): de arylgroep
wordt aan het eindstandige koolstofatoom gebonden, en de syn-$\beta$-hydride-eliminatie vanuit de stabielste conformeer
geeft het alkeen \textit{E}.
\end{solution}

\begin{solution}{exo:b3:homogeneous-catalysis:5}
De oxidatieve additie is traag en $\ce{[Rh(CO)2I2]-}$ is de rusttoestand, dus
$v = k[\mathrm{Rh}]_{\mathrm{tot}}[\ce{CH3I}]$. Methanol regenereert alleen methyljodide door een snelle reactie met HI;
het totaal ingebrachte jodide legt $[\ce{CH3I}]$ vast, zodat de snelheid niet van de
methanolconcentratie afhangt tot er te weinig methanol over is om HI weer
in methyljodide om te zetten.
\end{solution}

\begin{solution}{exo:b3:homogeneous-catalysis:6}
Fenylboorzuur met 4-broomtolueen, of 4-methylfenylboorzuur met
broombenzeen. Beide werken; de eerste gebruikt het goedkoopste boorzuur
en een stabiel, gangbaar arylbromide.
\end{solution}

\begin{solution}{exo:b3:homogeneous-catalysis:7}
Symmetrie $C_2$: isotactisch (beide plaatsen kiezen hetzelfde vlak). Symmetrie $C_s$:
syndiotactisch (spiegelbeeldige plaatsen wisselen af). Onoverbrugd $\ce{Cp2ZrCl2}$: atactisch
(achirale plaatsen).
\end{solution}

\begin{solution}{exo:b3:homogeneous-catalysis:8}
$\mathrm{TON} = 0.98/0.0010 = 980$; gemiddelde $\mathrm{TOF} = 980/\qty{2.0}{h} = \qty{490}{h^{-1}}$.
\end{solution}

\begin{solution}{exo:b3:homogeneous-catalysis:9}
De herhalingseenheid is $\ce{-[CH=CH-C5H8]-}$ met de twee CH-groepen op de koolstofatomen 1 en 3 van een
cyclopentaanring (poly(1,3-cyclopentyleenvinyleen)). Het vrijkomen van de
spanning van de bicyclische ring maakt $\Delta_rH^\circ$ sterk negatief, wat zwaarder weegt
dan het verlies aan translatie-entropie; er ontstaat geen gas, dus de
drijvende kracht is enthalpisch.
\end{solution}

\begin{solution}{exo:b3:homogeneous-catalysis:10}
$\bar n_n = 1/(1 - p) = 2000$; $M_w/M_n = 1 + p = 1.9995 \approx 2.0$. Een \omterm{def:b3:homogeneous-catalysis:ziegler-natta}{ziegler-nattakatalysator} heeft
meerdere soorten plaatsen, elk met haar eigen $p$; de som van meerdere
verdelingen van Schulz-Flory met verschillende gemiddelden is breder
(dispersiteit vaak 4 tot 8).
\end{solution}

\begin{solution}{exo:b3:homogeneous-catalysis:11}
Bij lage druk is $K_1K_2[\ce{H2}] \ll [\mathrm L] + K_1$: $v \approx kK_1K_2[\mathrm{Rh}]_{\mathrm{tot}}[\text{alkeen}][\ce{H2}]/([\mathrm L] + K_1)$,
van de eerste orde in $\ce{H2}$. Bij hoge druk overheerst de term $K_1K_2[\ce{H2}]$: $v \to k[\mathrm{Rh}]_{\mathrm{tot}}[\text{alkeen}]$,
onafhankelijk van $\ce{H2}$. $[\mathrm L]$ komt alleen in de noemer voor: toegevoegd fosfine verlaagt $v$
door het rodium terug te duwen naar de verzadigde prekatalysator.
\end{solution}

\begin{solution}{exo:b3:homogeneous-catalysis:12}
De insertie die rodium op het inwendige koolstofatoom zet, bouwt een
secundair alkyl op naast een volgepakt metaal; omvangrijke fosfines (en
een overmaat ervan, die er twee op het metaal houdt) maken die
overgangstoestand ongunstiger dan die welke het primaire alkyl geeft, dat
tot het onvertakte aldehyde leidt. Onvertakt butanal is het aldehyde dat
door aldolcondensatie en hydrogenering omgezet wordt in de vertakte alcohol $C_8$ die in
weekmakers gebruikt wordt; 2-methylpropanal is een minder waardevol
bijproduct.
\end{solution}

\begin{solution}{pb:b3:homogeneous-catalysis:1}
\textbf{1.} Ir(I), $d^8$: $9 + 4 + 2 + 1 = 16$ elektronen (neutrale telling, met de lading).

\textbf{2.} $\ce{[CH3Ir(CO)2I3]-}$, 18 elektronen, Ir(III).

\textbf{3.} $\ce{[CH3Ir(CO)3I2]}$, 18 elektronen, Ir(III).

\textbf{4.} $\ce{[CH3COIr(CO)2I2]}$, 16 elektronen, Ir(III).

\textbf{5.} $\ce{[CH3COIr(CO)2I3]-}$ (18) elimineert $\ce{CH3COI}$, wat $\ce{[Ir(CO)2I2]-}$ regenereert.

\textbf{6.} $\ce{CH3COI + H2O -> CH3COOH + HI}$; $\ce{CH3OH + HI -> CH3I + H2O}$.

\textbf{7.} $\ce{CH3OH + CO -> CH3COOH}$.

\textbf{8.} Met iridium is de migratorische insertie traag; de oxidatieve
additie van methyljodide, snelheidsbepalend met rodium, verloopt veel
sneller op het elektronenrijkere iridium.

\textbf{9.} Ze nemen één jodide weg van $\ce{[CH3Ir(CO)2I3]-}$, wat het neutrale tricarbonyl geeft, waarin
het metaal minder terugdoneert en de methylgroep sneller migreert naar
een elektrofieler CO.

\textbf{10.} Een omgekeerde afhankelijkheid van de concentratie vrij
jodide: jodide duwt het voorevenwicht terug naar het trage anionische
complex.

\textbf{11.} Methanol komt pas na de snelheidsbepalende stap binnen, en
wordt door een snelle reactie met HI omgezet in methyljodide.

\textbf{12.} $M(\ce{CH3COOH}) = \qty{60.1}{g/mol}$: $n = \qty{5.0e11}{g}/\qty{60.1}{g/mol} = \qty{8.3e9}{mol}$ per jaar.

\textbf{13.} $\qty{8.32e9}{mol}/\qty{8000}{h} = \qty{1.04e6}{mol/h}$.

\textbf{14.} $\qty{2.0e5}{L} \times \qty{5.0e-3}{mol/L} = \qty{1.0e3}{mol}$ iridium, $\qty{1.0e3}{mol} \times \qty{192.2}{g/mol} \approx \qty{190}{kg}$.

\textbf{15.} $\mathrm{TOF} = \qty{1.04e6}{mol/h}/\qty{1.0e3}{mol} = \qty{1.0e3}{h^{-1}}$, ongeveer \qty{0.29}{s^{-1}}.

\textbf{16.} $\mathrm{TON} = \qty{8.3e9}{mol}/\qty{1.0e3}{mol} = \num{8.3e6}$ per jaar.

\textbf{17.} $\qty{5.0e8}{kg}/(\qty{8000}{h} \times \qty{200}{m^3}) \approx \qty{310}{kg.m^{-3}.h^{-1}}$.

\textbf{18.} $\ce{CO + H2O -> CO2 + H2}$.

\textbf{19.} $1/0.94 = \qty{1.064}{mol}$ CO per mol azijnzuur.

\textbf{20.} $1.064 - 1 = \qty{0.064}{mol}$ \ce{CO2} per mol azijnzuur.

\textbf{21.} Eén ton is $\qty{1.0e6}{g}/\qty{60.1}{g/mol} = \qty{1.66e4}{mol}$; $\qty{1.66e4}{mol} \times 0.0638 \times \qty{44.0}{g/mol} \approx \qty{47}{kg}$ \ce{CO2}.

\textbf{22.} Water is een beginstof van de shiftreactie: minder water
vertraagt haar en bespaart de energie voor het drogen van het zuur. Water
blijft nodig om acetyljodide te hydrolyseren en om de katalysator actief
en opgelost te houden; met rodium slaat bij weinig water rodiumjodide
neer.

\textbf{23.} Waterstof: het hoopt zich op in het gas, verlaagt de
partiële druk van CO (zodat er gas afgeblazen moet worden, met verlies
van CO), en kan intermediairen hydrogeneren tot bijproducten zoals
propionzuur.

\textbf{24.} De iridiumkatalysator haalt ongeveer $\qty{1.0e3}{h^{-1}}$: elk iridiumatoom maakt
duizend moleculen azijnzuur per uur, zo'n acht miljoen per jaar.
\end{solution}
